\tocsection{\textbf{11/6/26 Lecture 29:} Continuity, Bases, and Subspaces}

\begin{definition}[\textbf{Definition 29.1:} Continuity]
    Let $(X, \cal{T})$, $(Y, \cal{U})$ be topological spaces. A map $f : X \to Y$ is \emph{continuous} if $f^{-1}(U) \in \cal{T}$ for all $U \in \cal{U}$.
\end{definition}

\begin{definition}[\textbf{Definition 29.2:} Convergence]
    Let $(X, \cal{T})$ be a topological space. Then, $x_n$ \emph{converges} to $x$ if every open set $U \ni x$ contains all but finitely many points of the sequence.
\end{definition}

\begin{lemma}[\textbf{Proposition 29.3}]
    Let $f : X \to Y$ be continuous. Then, if $x_n \to x$, then $f(x_n) \to f(x)$. Furthermore, if $X$ is a metric space and for all $x$, $x_n \to x \implies f(x_n) \to f(x)$, then $f$ is continuous.
\end{lemma}
\begin{proof}{Proof:}
    Let $U$ be an open set containing $f(x)$. $f^{-1}(U)$ is an open set containing $x$, so it contains all but finitely many of the $x_n$. Then, all but finitely many of the $f(x_n)$ are contained in $U$.\\

    Now, suppose that $f$ is not continuous (in the topological sense). Then, there exists an open $U$ such that $f^{-1}(U)$ is not open. Thus, there exists $x \in f^{-1}(U)$ such that $B_r(x) \not\subseteq f^{-1}(U)$ for all $r > 0$. Picking $x_n$ so that $x_n \in B_{1/n}(x) \setminus f^{-1}(U)$, we see that $x_n \to x$, but as $f(x_n) \notin U$ for all $n$, $f(x_n) \not\to f(x)$, which is a contradiction.

    $\qed$
\end{proof}

\begin{example}{\textbf{Example:} The Co-Countable Topology}
    Consider $f : (\R, \cal{T}) \to (\R, \abs{\cdot})$ by $f(x) = x$, where $(\R, \abs{\cdot})$ has the usual metric topology. Let
    \[ \cal{T} = \{U \subseteq \R \mid U^{c} \tt{ is countable}\} \cup \{\emptyset\} \]
    We need to show $x_n \to x \implies f(x_n) \to f(x)$. Suppose $x_n \to x$. Let $A = \{x_n \mid x_n \neq x\}$. Now, $A^{c}$ is open and contains $x$, which implies $A^{c}$ contains all but finitely many of the $x_n$. Thus, for all but finitely many $x_n$, $x_n = x$.\\

    Thus, $f(x_n) \to f(x)$, since every open set containing $f(x)$ contains all $f(x_n)$ with $x_n = x$.\\

    However, $f$ is not continuous, since the preimage of nearly every open set is not open in $\cal{T}$. Thus, the metric topology requirement in the proposition above is important.
\end{example}

\begin{definition}[\textbf{Definition 29.4:} Sub-Bases and Bases]
    Let $X$ be a set and $\cal{B} \subseteq \P(X)$. We say that $\cal{B}$ is a \emph{sub-basis} if every $x \in X$ is contained in some $B \in \cal{B}$. We say a sub-basis $\cal{B}$ is a \emph{basis} if whenever $x \in B_1 \cap B_2$ for $B_1, B_2 \in \cal{B}$, there exists $B_3 \in \cal{B}$ such that $x \in B_3 \subseteq B_1 \cap B_2$.
\end{definition}

\begin{lemma}[\textbf{Proposition 29.5}]
    Let $X$ be a set and $\cal{B}$ a sub-basis.
    \begin{enumerate}
        \item Let $\cal{T}$ denote the collection of subsets of $X$ of the form
        \[ O = \bigcup U_{\alpha} \]
        where $U_{\alpha}$ is the intersection of finitely many elements of $\cal{B}$. Then, $\cal{T}$ is a topology.
        \item If $\cal{B}$ is a basis, and $\cal{U}$ is the collection of arbitrary unions of sets in $\cal{B}$, then $\cal{U}$ is a topology.
    \end{enumerate}
\end{lemma}
\begin{proof}{Proof:}
    \textbf{(1)} Since $\cal{B}$ is a sub-basis, $X \in \cal{T}$, and taking an empty union, $\emptyset \in \cal{T}$. Arbitrary unions are immediate. Then, for
    \[ O_1 = \bigcup U_{\alpha} \qquad O_2 = \bigcup U_{\beta} \]
    we see
    \[ O_1 \cap O_2 = \bigcup_{\alpha, \beta} U_{\alpha} \cap U_{\beta} \]
    which is of the desired form. Thus, $\cal{T}$ is a topology.\\

    \textbf{(2)} As before, we just need to check finite intersections. Let $O_1, O_2 \in \cal{U}$ and let $x \in O_1 \cap O_2$. Then, $x \in B_1 \subseteq O_1$ and $x \in B_2 \subseteq O_2$. Thus, $x \in B_x \subseteq B_1 \cap B_2 \subseteq O_1 \cap O_2$. Then, we can write
    \[ O_1 \cap O_2 = \bigcup_{x \in O_1 \cap O_2} B_x \]
    which is in $\cal{U}$.

    $\qed$
\end{proof}

We call the topologies above the \emph{topology generated by $\cal{B}$}.

\begin{example}{\textbf{Example:} Countable Bases for $\R^2$}
    A countable basis for $\R^2$ with the metric topology is given by
    \[ \{(a, b) \times (c, d) \col a, b, c, d \in \Q\} \]
    An example of a countable sub-basis for $\R^2$ is given by
    \[ \{(a, b) \times \R \col a, b \in \Q\} \cup \{\R \times (c, d) \col c, d \in \Q\} \]
\end{example}

Given a topological space $(X, \cal{T})$ and $Y \subseteq X$, we can give this a topology by letting $Z \subseteq Y$ be open if and only if it is given by $U \cap Y$ for $U \in \cal{T}$. We call this the \emph{subspace topology} on $Y$.

\begin{example}{\textbf{Example:} The Circle}
    Let $X = \R^2$ and $S^1 = \{x \in \R^2 \col \abs{x} = 1\}$.
    \fig{0.7}{Lecture29Fig1}{}
    Note that the arc alone is not open in $\R^2$, but is in $S^1$.
\end{example}

\begin{definition}[\textbf{Definition 29.6:} Homeomorphisms]
    A map $f : X \to Y$ is a \emph{homeomorphism} if it is a continuous bijection and its inverse is continuous.
\end{definition}

\begin{example}{\textbf{Example:} $1/x$}
    The map $f : (0, 1) \to (1, \infty)$ by $f(x) = \frac{1}{x}$ is a homeomorphism.
\end{example}
